% !PRINTED 467
\weberpage{467}

%% [running head, not body text -- commented out 2026-08-07] Page furniture (running
%% head, folio, signature) is excluded from this edition by standing policy.
%% \noindent \S\,127. \hfill The powers of $f(\omega)$ as class invariants. \hfill 467

\S\,127. \textbf{\gesperrt{The Powers of $f(\omega)$ as Class Invariants.}}

By the definition of $\gamma_2(\omega)$ we have:
\weq{1}{f(\omega)^8=\dfrac{f(\omega)^{24}-16}{\gamma_2(\omega)}.}
If $\omega$ is the root of an equation:
\weq{2}{A\omega^2+2B\omega+C=0,\quad AC-B^2=m,}
whose discriminant $-4m$ is not divisible by 3, then we may assume, passing if necessary to an equivalent equation, that $A$, $C$ are odd, that $A$ is indivisible by 3 and that $B$ is divisible by 3. But then, by the two preceding sections, $f(\omega)^{24}$ and $\gamma_2(\omega)$ are class invariants, and from (1) we obtain the theorem:

3. \gesperrt{If $\omega$ is the root of a quadratic form of negative discriminant not divisible by 3, whose two outer coefficients are odd and whose middle coefficient is divisible by 6, then}
\[
f(\omega)^8
\]
\gesperrt{is a class invariant.}

In order to investigate the question whether still lower powers of $f(\omega)$ can also be class invariants, we make use of the formula:
\weq{3}{f(\omega+2r)^6=i^{-r}f(\omega)^6,}
in which $r$ is an integer; the values
\[
\omega_r=\omega+2r
\]
are the roots of equivalent quadratic equations
\[
A\omega_r^2+2B_r\omega_r+C_r=0,
\]
in which
\weq{4}{B_r=B-2Ar,}
and here $r$ can be determined according to the modulus 4 in such a way that, for odd $m$:
\weq{5}{B_r\equiv 0,2,4,6\pmod{8}}
and for even $m$:
\weq{6}{B_r\equiv 1,3,5,7\pmod{8}.}
If we now succeed in finding an equation with rational coefficients that, among the four values (3), is satisfied either only
%CONT>>
